% ============================================================
%  KICH BAN VIDEO CA NHAN --- DAT
%  Vai tro: Leader / Security Engineer
%  Files: privacy.py | anonymizer_backend.py | clip_recorder.py | admin_auth.py
%  Thoi luong muc tieu: 13-15 phut
% ============================================================
\documentclass[12pt, a4paper]{article}
\usepackage[utf8]{inputenc}
\usepackage[T5]{fontenc}
\usepackage[vietnamese]{babel}
\usepackage[margin=2cm, top=2.5cm, bottom=2.5cm]{geometry}
\usepackage{parskip}
\usepackage[dvipsnames]{xcolor}
\usepackage{tcolorbox}
\tcbuselibrary{skins, breakable, listings}
\usepackage{listings}
\usepackage{booktabs}
\usepackage{tabularx}
\usepackage{fontawesome5}
\usepackage{pifont}
\newcommand{\cmark}{\ding{51}}
\newcommand{\xmark}{\ding{55}}

\lstset{
  basicstyle=\ttfamily\footnotesize,
  keywordstyle=\color{NavyBlue}\bfseries,
  commentstyle=\color{OliveGreen}\itshape,
  stringstyle=\color{BrickRed},
  numbers=left, numberstyle=\tiny\color{gray}, numbersep=6pt,
  frame=single, breaklines=true, showstringspaces=false,
  tabsize=4, backgroundcolor=\color{gray!7},
}

\usepackage{titlesec}
\titleformat{\section}[block]
  {\large\bfseries\color{BrickRed}}{\thesection.}{0.5em}{}[\titlerule]
\titleformat{\subsection}[block]
  {\normalsize\bfseries\color{BrickRed!80}}{\thesubsection}{0.5em}{}

\usepackage{fancyhdr}
\pagestyle{fancy}\fancyhf{}
\lhead{\textcolor{BrickRed}{\textbf{Kịch Bản Video --- ĐẠT (Leader / Security)}}}
\rhead{\textcolor{gray}{XLA Face Privacy System}}
\cfoot{\thepage}
\usepackage[hidelinks]{hyperref}

%% ---- Environments ----
\newtcolorbox{say}[1][]{
  title={\faMicrophone\ \textbf{NÓI}},
  colback=BrickRed!4, colframe=BrickRed!70,
  fonttitle=\small\bfseries, left=4pt, right=4pt, breakable, #1}

\newtcolorbox{doact}[1][]{
  title={\faMousePointer\ \textbf{THAO TÁC}},
  colback=OliveGreen!5, colframe=OliveGreen!70,
  fonttitle=\small\bfseries, left=4pt, right=4pt, breakable, #1}

\newtcolorbox{cam}[1][]{
  title={\faVideo\ \textbf{CAMERA}},
  colback=Goldenrod!8, colframe=Goldenrod!80,
  fonttitle=\small\bfseries, left=4pt, right=4pt, breakable, #1}

\newtcolorbox{code}[2][]{
  title={\faCode\ \textbf{CODE --- #2}},
  colback=NavyBlue!4, colframe=NavyBlue!55,
  fonttitle=\small\bfseries\color{NavyBlue}, left=4pt, right=4pt, breakable, #1}

\newtcolorbox{warn}[1][]{
  title={\faExclamationTriangle\ \textbf{CHÚ Ý}},
  colback=red!4, colframe=red!60,
  fonttitle=\small\bfseries, left=4pt, right=4pt, breakable, #1}

\newtcolorbox{timing}[1]{
  colback=NavyBlue!10, colframe=NavyBlue!50,
  title={\faClock\ \textbf{#1}},
  fonttitle=\small\bfseries\color{NavyBlue}, left=4pt, right=4pt}

\newcommand{\fpath}[1]{\texttt{\textcolor{purple}{#1}}}
\newcommand{\lnum}[1]{\colorbox{yellow!40}{\texttt{\small dòng~#1}}}
\newcommand{\pause}{\textcolor{gray}{\textit{[dừng 1 giây]}}}
\newcommand{\camlook}{\textcolor{Goldenrod}{\textbf{[nhìn camera]}}}
\newcommand{\screen}{\textcolor{OliveGreen}{\textbf{[nhìn màn hình]}}}

% ============================================================
\begin{document}
% ============================================================

\begin{titlepage}
\centering\vspace*{1.5cm}
{\Huge\bfseries\color{BrickRed} KỊCH BẢN VIDEO\par}
\vspace{0.3cm}
{\LARGE\bfseries\color{BrickRed!80} ĐẠT --- Leader / Security Engineer\par}
\vspace{0.5cm}
{\large\color{gray} XLA Face Privacy System --- Đồ Án Cuối Kỳ\par}
\vspace{1.2cm}

\begin{tcolorbox}[colback=BrickRed!7, colframe=BrickRed, width=0.9\textwidth]
\centering\large
\textbf{Em trực tiếp viết 4 files:}\\[0.4cm]
\renewcommand{\arraystretch}{1.5}
\begin{tabular}{llr}
  \hline
  \fpath{src/privacy.py} & Blur Engine — 7 mode ẩn danh & \textbf{269 dòng} \\
  \fpath{src/anonymizer\_backend.py} & Config \& 4 Preset & \textbf{237 dòng} \\
  \fpath{src/clip\_recorder.py} & AES-256-GCM Encrypt/Decrypt & \textbf{397 dòng} \\
  \fpath{src/admin\_auth.py} & PBKDF2 + Timing Attack Prevention & \textbf{130 dòng} \\
  \hline
  \multicolumn{2}{l}{\textbf{Tổng cộng}} & \textbf{1033 dòng} \\
  \hline
\end{tabular}
\end{tcolorbox}

\vspace{0.8cm}
\begin{tcolorbox}[colback=gray!6, colframe=gray!50, width=0.9\textwidth]
\textbf{Thời lượng mục tiêu:} 13--15 phút\quad
\textbf{Định dạng:} Quay một mình, camera rõ mặt\\[0.2cm]
\textbf{Cần chuẩn bị:}
Server chạy cổng 8000 \textbullet{}
VS Code font $\geq$16 \textbullet{}
4 files đã mở sẵn \textbullet{}
Browser tab dashboard và secure-archive \textbullet{}
Tắt notification Windows
\end{tcolorbox}

\vspace{0.8cm}
\begin{warn}
Trước khi quay: đọc lại toàn bộ script này ÍT NHẤT 2 lần.
Luyện nói to trước gương hoặc tự record thử 2 phút đầu.
Không cần học thuộc --- nhìn script và nói tự nhiên là được.
\end{warn}
\end{titlepage}

\tableofcontents\newpage

% ============================================================
\section{CHUẨN BỊ TRƯỚC KHI BẤM REC}
% ============================================================

\begin{doact}
\textbf{Bước 1 --- Khởi động server (terminal):}
\begin{lstlisting}[language=bash, numbers=none]
cd C:\Users\itent\Desktop\xla_demo
.venv\Scripts\activate
uvicorn src.app_api:app --host 0.0.0.0 --port 8000 --reload
\end{lstlisting}
Chờ thấy \texttt{Application startup complete.} mới tiếp tục.
\end{doact}

\begin{doact}
\textbf{Bước 2 --- Mở VS Code, mở sẵn 4 tab files:}\\
Tab 1: \fpath{src/privacy.py} (269 dòng)\\
Tab 2: \fpath{src/anonymizer\_backend.py} (237 dòng)\\
Tab 3: \fpath{src/clip\_recorder.py} (397 dòng)\\
Tab 4: \fpath{src/admin\_auth.py} (130 dòng)\\
Font size: \textbf{16--18}. Zoom màn hình 125\%.
\end{doact}

\begin{doact}
\textbf{Bước 3 --- Mở 2 tab browser:}\\
Tab 1: \texttt{http://localhost:8000/admin/dashboard.html}\\
Tab 2: \texttt{http://localhost:8000/admin/secure-archive.html}
\end{doact}

\begin{cam}
Camera đặt ngang tầm mắt, cách mặt 60--80cm.
Ánh sáng từ phía trước mặt (không ngồi trước cửa sổ).
Đặt màn hình bên phải hoặc trái camera để vừa nói vừa có thể nhìn script.
\end{cam}

\begin{warn}
Test âm thanh 10 giây trước khi quay.
Nháy mắt bình thường --- không nhìn thẳng camera quá căng.
Nói chậm hơn bình thường khoảng 20\% --- video thường nghe nhanh hơn thực tế.
\end{warn}

\newpage
% ============================================================
\section{0:00--0:15 --- MỞ ĐẦU}
% ============================================================

\begin{timing}{0:00 -- 0:15}
Ngồi thẳng, nhìn camera. Thở sâu. Bắt đầu sau khi đã ổn định.
\end{timing}

\begin{cam}
Nhìn thẳng vào camera. Giữ ánh mắt tự nhiên. Nói từ tốn.
\end{cam}

\begin{say}
``Xin chào thầy cô. Em là Đạt.

Trong đồ án này em đảm nhận vai trò Leader nhóm
và phụ trách toàn bộ phần \textbf{bảo mật} của hệ thống.

Video này em sẽ trình bày \textbf{trực tiếp những đoạn code em đã viết}---
chỉ rõ từng dòng, tại sao viết như vậy, và quyết định kỹ thuật nào đằng sau nó.''
\end{say}

% ============================================================
\section{0:15--1:30 --- GIỚI THIỆU ĐÓNG GÓP CỤ THỂ}
% ============================================================

\begin{timing}{0:15 -- 1:30}
Chuyển góc nhìn sang màn hình, mở VS Code.
\end{timing}

\begin{doact}
Mở VS Code. Hiện cây thư mục bên trái. Hover chuột lần lượt qua từng file khi nói tên.
\end{doact}

\begin{say}
``Phần em trực tiếp viết gồm bốn files.

File thứ nhất: \textbf{privacy.py} ở thư mục src --- 269 dòng.
Đây là toàn bộ Blur Engine: 7 mode ẩn danh khuôn mặt, từ pixelate đơn giản đến obliterate hoàn toàn không thể nhận ra.

File thứ hai: \textbf{anonymizer\_backend.py} --- 237 dòng.
Config system và 4 preset: fast, balanced, strong, strict.

File thứ ba: \textbf{clip\_recorder.py} --- 397 dòng.
Mã hoá và giải mã clip với AES-256-GCM --- tiêu chuẩn mã hoá công nghiệp.

File thứ tư: \textbf{admin\_auth.py} --- 130 dòng.
Xác thực admin với PBKDF2 và phòng chống timing attack.

Tổng cộng: 1033 dòng code em tự viết.''

\pause\ \camlook\ ``Em bắt đầu với phần quan trọng nhất: Blur Engine.''
\end{say}

% ============================================================
\section{1:30--2:00 --- ĐẶT VẤN ĐỀ: TẠI SAO CẦN BLUR ENGINE}
% ============================================================

\begin{timing}{1:30 -- 2:00}
Nhìn camera, không cần mở code ngay.
\end{timing}

\begin{say}
``Trước khi vào code, em muốn giải thích bài toán.

Camera giám sát quay tất cả mọi người---kể cả người không đồng ý lưu khuôn mặt.
Nguyên tắc \textbf{Privacy by Design} theo GDPR Article 25 yêu cầu:
ẩn danh hoá phải là mặc định, không phải tuỳ chọn thêm.

Nhưng nếu chỉ lưu bản đã blur, bằng chứng điều tra bị mất.

Giải pháp em thiết kế: \textbf{Dual Storage} ---
lưu song song bản blur cho người dùng thông thường
và bản gốc mã hoá AES-256-GCM chỉ admin mới mở được.

Blur Engine phải đủ mạnh để không thể nhận ra khuôn mặt,
nhưng cũng đủ nhanh để chạy real-time 25--30 FPS.''
\end{say}

% ============================================================
\section{2:00--6:00 --- PRIVACY.PY --- BLUR ENGINE}
% ============================================================

\begin{timing}{2:00 -- 6:00}
Chuyển sang VS Code, mở tab \fpath{src/privacy.py}.
\end{timing}

\subsection{2:00--2:30 --- Nhìn tổng quan file}

\begin{doact}
Nhấn \texttt{Ctrl+Shift+E} mở Explorer, click vào \fpath{src/privacy.py}.
Scroll từ đầu xuống cuối nhanh trong 5 giây để thầy cô thấy độ dài (269 dòng).
Sau đó nhấn \texttt{Ctrl+Home} về đầu file.
\end{doact}

\begin{say}
``Đây là \fpath{privacy.py} --- 269 dòng em viết hoàn toàn.

Nhìn tổng quan: dòng 1--10 là import và helper nhỏ.
Dòng 67 trở đi là 7 hàm blur độc lập.
Dòng 201 là hàm \texttt{anonymize\_faces} tổng hợp --- pipeline gọi vào đây.

Em sẽ đi qua từng hàm quan trọng.''
\end{say}

\subsection{2:30--3:30 --- Hàm \_fast\_gaussian\_blur (dòng 67--76)}

\begin{doact}
Nhấn \texttt{Ctrl+G}, gõ \textbf{67}, Enter. Highlight từ dòng 67 đến 76.
\end{doact}

\begin{say}
``Dòng 67: hàm \texttt{\_fast\_gaussian\_blur}.
Đây là blur mode cơ bản, được dùng trong tất cả các preset làm bước pre-processing.''
\end{say}

\begin{code}{privacy.py dòng 67--76 --- \_fast\_gaussian\_blur()}
\begin{lstlisting}[language=Python, firstnumber=67]
def _fast_gaussian_blur(roi: np.ndarray,
                        blur_scale: float,
                        blur_kernel: int) -> np.ndarray:
    h, w = roi.shape[:2]
    scale = min(max(blur_scale, 0.05), 1.0)
    down_w = max(1, int(round(w * scale)))
    down_h = max(1, int(round(h * scale)))
    small = cv2.resize(roi, (down_w, down_h),
                       interpolation=cv2.INTER_LINEAR)
    small = cv2.GaussianBlur(small,
                 (_odd(blur_kernel), _odd(blur_kernel)), 0)
    return cv2.resize(small, (w, h),
                      interpolation=cv2.INTER_LINEAR)  # dong 76
\end{lstlisting}
\end{code}

\begin{say}
\screen\ ``Dòng 68 nhận 3 tham số: ROI là vùng khuôn mặt đã cắt ra, blur\_scale và kernel size.

Dòng 71--72: \textbf{downscale ROI} xuống còn \texttt{scale\%} kích thước.
Preset balanced dùng scale $\approx$20\% --- mặt 200×200px shrink xuống còn 40×40px.

Dòng 73--74: GaussianBlur trên ảnh \textit{nhỏ} --- \textbf{đây là kỹ thuật tối ưu tốc độ}.
Blur trực tiếp trên 200×200px: $200 \times 200 = 40.000$ pixels.
Blur qua downscale: chỉ $40 \times 40 = 1.600$ pixels.
Nhanh hơn 25 lần với kết quả hình ảnh gần như tương đương.

Dòng 75--76: upscale lại, kết quả là blur mềm mại, tự nhiên.''
\end{say}

\subsection{3:30--4:15 --- Hàm \_pixelate (dòng 78--86)}

\begin{doact}
Nhấn \texttt{Ctrl+G}, gõ \textbf{78}, Enter. Highlight dòng 78--86.
\end{doact}

\begin{say}
``Dòng 78: \texttt{\_pixelate} --- tạo hiệu ứng mosaic vuông.''
\end{say}

\begin{code}{privacy.py dòng 78--86 --- \_pixelate()}
\begin{lstlisting}[language=Python, firstnumber=78]
def _pixelate(roi: np.ndarray, pixel_block: int) -> np.ndarray:
    h, w = roi.shape[:2]
    block = max(2, pixel_block)
    small_w = max(1, w // block)
    small_h = max(1, h // block)
    small = cv2.resize(roi, (small_w, small_h),
                       interpolation=cv2.INTER_LINEAR)
    return cv2.resize(small, (w, h),
                      interpolation=cv2.INTER_NEAREST)  # dong 86
\end{lstlisting}
\end{code}

\begin{say}
``Dòng 84: downscale xuống 1/block kích thước.

\textbf{Chìa khoá nằm ở dòng 86}: \texttt{INTER\_NEAREST} thay vì INTER\_LINEAR.
INTER\_NEAREST không nội suy --- copy pixel gần nhất.
Kết quả: mỗi block vuông y hệt nhau, tạo mosaic đặc trưng như ảnh game 8-bit.
Nếu dùng INTER\_LINEAR sẽ ra blur, không ra pixelate.

Preset \textbf{fast} (\fpath{anonymizer\_backend.py} dòng 55) dùng \texttt{pixel\_block=18}
--- mỗi ô mosaic 18×18px.''
\end{say}

\subsection{4:15--6:00 --- Hàm \_obliterate\_features (dòng 172--200) --- QUAN TRỌNG NHẤT}

\begin{doact}
Nhấn \texttt{Ctrl+G}, gõ \textbf{172}, Enter. Expand cửa sổ code rộng ra.
\end{doact}

\begin{say}
``Đây là hàm quan trọng nhất em viết trong file này.
\texttt{\_obliterate\_features} --- mode \textbf{strong} của hệ thống.

Mục tiêu: không cần blur mượt mà --- cần \textbf{phá huỷ hoàn toàn}.
Kể cả các thuật toán AI nhận diện mặt tốt nhất cũng không thể phục hồi.''
\end{say}

\begin{code}{privacy.py dòng 172--200 --- \_obliterate\_features() --- 4 bước phá huỷ}
\begin{lstlisting}[language=Python, firstnumber=172]
def _obliterate_features(
    roi: np.ndarray,
    obliterate_scale: float,  # preset strong: 0.06
    blur_kernel: int,
    scramble_block: int,      # preset strong: 10
    palette_levels: int,      # preset strong: 8
    seed: int,
) -> np.ndarray:
    h, w = roi.shape[:2]
    scale = min(max(obliterate_scale, 0.03), 0.25)
    down_w = max(1, int(round(w * scale)))
    down_h = max(1, int(round(h * scale)))

    # BUOC 1 (dong 185): compress xuong 6%, xoa toan bo dac diem
    work = cv2.resize(roi, (down_w, down_h),
                      interpolation=cv2.INTER_AREA)
    work = cv2.resize(work, (w, h),
                      interpolation=cv2.INTER_NEAREST)

    # BUOC 2 (dong 191): scramble block 10x10
    work = _scramble_blocks(work, block_size=scramble_block, seed=seed)

    # BUOC 3 (dong 194): GaussianBlur kernel >= 35px
    work = cv2.GaussianBlur(work,
               (_odd(max(blur_kernel, 35)),
                _odd(max(blur_kernel, 35))), 0)

    # BUOC 4 (dong 198): quantize ve 8 mau
    work = _quantize_colors(work, levels=palette_levels)
    return work  # dong 200
\end{lstlisting}
\end{code}

\begin{say}
``Hàm nhận preset strong với các tham số: \texttt{obliterate\_scale=0.06},
\texttt{scramble\_block=10}, \texttt{palette\_levels=8}.

\textbf{Bước 1 --- dòng 185--188: Nén cực mạnh.}
Compress khuôn mặt xuống còn \textbf{6\% kích thước}.
Ví dụ mặt 200×200px $\to$ 12×12px $\to$ upscale lại 200×200px.
Kết quả: mọi đặc điểm nhận dạng --- mắt, mũi, miệng, hình dạng mặt --- biến mất.
Không còn gì để AI nhận diện.

\textbf{Bước 2 --- dòng 191: Xáo trộn không gian.}
\texttt{\_scramble\_blocks} xáo trộn các block 10×10px theo thứ tự ngẫu nhiên (seed cố định).
Nếu bước 1 vẫn còn sót lại vài điểm texture, bước 2 phá vỡ quan hệ không gian.
Phần trán có thể nằm ở vị trí cằm --- không thể reconstruct.

\textbf{Bước 3 --- dòng 194--197: Blur cứng.}
GaussianBlur kernel tối thiểu 35×35px --- lớn hơn nhiều so với normal blur.
Xoá những cạnh sắc sót lại sau bước 1 và 2.

\textbf{Bước 4 --- dòng 198: Giảm màu sắc.}
Quantize xuống còn 8 màu. Loại bỏ micro-texture và gradient tinh tế.
Các mạng neural hiện đại có thể khai thác gradient dù rất nhỏ --- bước này chặn điều đó.

Bốn bước kết hợp: \textbf{không thể đảo ngược bằng bất kỳ phương pháp nào.}''
\end{say}

% ============================================================
\section{6:00--7:30 --- ANONYMIZER\_BACKEND.PY --- 4 PRESETS}
% ============================================================

\begin{timing}{6:00 -- 7:30}
Chuyển sang tab \fpath{src/anonymizer\_backend.py}, nhảy dòng 53.
\end{timing}

\begin{doact}
Click tab \fpath{anonymizer\_backend.py} trong VS Code.
Nhấn \texttt{Ctrl+G}, gõ \textbf{53}, Enter.
\end{doact}

\begin{say}
``File này em thiết kế hệ thống config cho toàn bộ pipeline.
Quan trọng nhất là dict PRESETS từ dòng 53.''
\end{say}

\begin{code}{anonymizer\_backend.py dòng 53--91 --- PRESETS dict --- 4 chế độ}
\begin{lstlisting}[language=Python, firstnumber=53]
PRESETS: Dict[PresetName, AnonymizationConfig] = {
    "fast": AnonymizationConfig(
        mode="pixelate",      # goi _pixelate() -> privacy.py dong 78
        imgsz=416,            # YOLO input nho -> nhanh hon
        detect_every=3,       # detect 1 trong 3 frame -> 3x tiet kiem CPU
        proc_width=512,
        pixel_block=18,
    ),
    "balanced": AnonymizationConfig(
        mode="headcloak",     # che toan bo vung dau
        imgsz=512,
        detect_every=2,       # detect 1/2 frame
        proc_width=640,
        head_ratio=0.65,      # vung che = 65% chieu cao dau
        face_padding=0.16,
    ),
    "strong": AnonymizationConfig(
        mode="obliterate",    # goi _obliterate_features() -> dong 172
        imgsz=512,
        detect_every=2,
        obliterate_scale=0.06,
        scramble_block=10,
        palette_levels=8,
    ),
    "strict": AnonymizationConfig(
        mode="headcloak",
        imgsz=416,
        detect_every=2,
        head_ratio=0.95,      # che 95% dau -> gan nhu toan bo
        face_padding=0.24,
        blur_kernel=39,       # kernel lon hon -> manh hon
    ),
}   # dong 91
\end{lstlisting}
\end{code}

\begin{say}
``Dòng 53: dict PRESETS với 4 key.

\textbf{fast} dòng 54--60: pixelate, YOLO imgsz=416 nhỏ hơn $\to$ nhanh hơn,
detect 1 trong 3 frame thay vì mỗi frame.
Phù hợp máy yếu hoặc cần FPS cao.

\textbf{balanced} dòng 62--68: headcloak --- che toàn bộ vùng đầu head\_ratio=0.65.
detect\_every=2: xử lý 15fps thay vì 30fps.
Cân bằng tốc độ và độ bảo mật.

\textbf{strong} dòng 70--78: mode obliterate --- gọi hàm dòng 172 vừa xem.
obliterate\_scale=0.06: đúng cái tham số em vừa giải thích là nén xuống 6\%.

\textbf{strict} dòng 80--86: headcloak với head\_ratio=0.95 --- che 95\% đầu,
blur\_kernel=39 lớn hơn nhiều.
Dùng khi yêu cầu bảo mật cao nhất, chấp nhận FPS thấp hơn.

Khi frontend gọi \texttt{GET /presets} (\fpath{app\_api.py} dòng 988),
Khánh lấy chính dict này trả về JSON.''
\end{say}

% ============================================================
\section{7:30--10:30 --- CLIP\_RECORDER.PY --- AES-256-GCM}
% ============================================================

\begin{timing}{7:30 -- 10:30}
Chuyển sang tab \fpath{src/clip\_recorder.py}.
\end{timing}

\begin{doact}
Click tab \fpath{clip\_recorder.py}.
Nhấn \texttt{Ctrl+G}, gõ \textbf{73}, Enter.
\end{doact}

\begin{say}
``File này em viết hệ thống mã hoá và giải mã clip video.
Bắt đầu từ dòng 73.''
\end{say}

\subsection{7:30--8:30 --- \_pack\_frames (dòng 73--86)}

\begin{code}{clip\_recorder.py dòng 73--86 --- \_pack\_frames(): serialize trước khi mã hoá}
\begin{lstlisting}[language=Python, firstnumber=73]
def _pack_frames(frames: List[np.ndarray], fps: float) -> bytes:
    """Serialise frames as length-prefixed JPEG stream."""
    buf = io.BytesIO()
    buf.write(struct.pack(">I", len(frames)))    # dong 76: 4B so frame
    buf.write(struct.pack(">I", int(fps * 1000))) # dong 77: 4B fps*1000
    for f in frames:
        ok, enc = cv2.imencode(
            ".jpg", f, [cv2.IMWRITE_JPEG_QUALITY, 90])  # dong 80: JPEG Q90
        if not ok:
            continue
        data = enc.tobytes()
        buf.write(struct.pack(">I", len(data)))  # dong 84: 4B do dai
        buf.write(data)                          # dong 85: JPEG bytes
    return buf.getvalue()                        # dong 86
\end{lstlisting}
\end{code}

\begin{say}
``\texttt{\_pack\_frames} là bước chuẩn bị: serialize danh sách NumPy frame thành bytes
trước khi đưa vào AES encrypt.

\textbf{Dòng 76--77}: header gồm 4 byte số frame và 4 byte fps nhân 1000.
Tại sao nhân 1000? Để tránh arithmetic float --- fps=29.97 thành 29970, integer chính xác.

\textbf{Dòng 80}: mỗi frame encode thành JPEG quality 90.
Quality 90 giữ đủ độ sắc nét cho bằng chứng, giảm size $\approx$5--8x so với raw.

\textbf{Dòng 84--85}: length-prefixed format --- mỗi frame ghi 4 byte kích thước trước.
Format này tự mô tả, không cần delimiter riêng.
Khi decrypt và đọc lại: đọc 4 byte header $\to$ biết frame tiếp theo dài bao nhiêu byte.''
\end{say}

\subsection{8:30--9:45 --- encrypt\_clip (dòng 109--119)}

\begin{doact}
Nhấn \texttt{Ctrl+G}, gõ \textbf{109}, Enter. Highlight dòng 109--119.
\end{doact}

\begin{code}{clip\_recorder.py dòng 109--119 --- encrypt\_clip() --- AES-256-GCM}
\begin{lstlisting}[language=Python, firstnumber=109]
def encrypt_clip(frames: List[np.ndarray],
                 fps: float,
                 password: str) -> bytes:
    """
    Returns: [ 32B salt || 12B nonce || ciphertext+AuthTag ]
    """
    plaintext = _pack_frames(frames, fps)   # dong 114: serialize frames
    salt  = secrets.token_bytes(32)         # dong 115: 256-bit random salt
    nonce = secrets.token_bytes(12)         # dong 116: 96-bit random nonce
    key   = derive_clip_key(password, salt) # dong 117: PBKDF2 -> 256-bit key
    ct    = AESGCM(key).encrypt(            # dong 118: AES-GCM encrypt
                nonce, plaintext, None)
    return salt + nonce + ct               # dong 119: self-contained blob
\end{lstlisting}
\end{code}

\begin{say}
``Đây là hàm quan trọng nhất của file này. Chỉ có 6 dòng nội dung nhưng mỗi dòng là một quyết định thiết kế.

\textbf{Dòng 114}: gọi \texttt{\_pack\_frames} vừa xem --- bytes thuần tuý.

\textbf{Dòng 115}: \texttt{secrets.token\_bytes(32)} --- 32 byte = 256 bit ngẫu nhiên.
Module \texttt{secrets} dùng CSPRNG của hệ điều hành, không phải \texttt{random.randbytes}.
Quan trọng: salt được tạo mới cho \textbf{mỗi clip}.
Nghĩa là: cùng password, clip A và clip B có key hoàn toàn khác nhau.
Hack được một clip không giúp hack các clip khác.

\textbf{Dòng 116}: nonce 12 byte = 96 bit.
NIST SP 800-38D khuyến nghị 96-bit là tối ưu cho AES-GCM.
Cũng random mới mỗi clip.

\textbf{Dòng 117}: \texttt{derive\_clip\_key} gọi sang \fpath{admin\_auth.py} dòng 124
--- đó là PBKDF2-SHA256 với 390.000 iterations.
Password ngắn $\to$ key 256-bit đủ mạnh.

\textbf{Dòng 118}: \texttt{AESGCM(key).encrypt} --- AES-256 chế độ GCM.
GCM = Galois/Counter Mode: vừa mã hoá vừa tạo Authentication Tag 128-bit.
Bất kỳ thay đổi nào trên ciphertext sẽ khiến tag sai $\to$ decrypt fail.
Đây là khác biệt lớn so với AES-CBC chỉ mã hoá, không có integrity check.

\textbf{Dòng 119}: file output tự chứa hết mọi thứ: salt + nonce + ciphertext.
Không cần bảng mapping ngoài, không cần metadata file riêng.''
\end{say}

\subsection{9:45--10:30 --- decrypt\_clip (dòng 122--138)}

\begin{doact}
Nhấn \texttt{Ctrl+G}, gõ \textbf{122}, Enter.
\end{doact}

\begin{code}{clip\_recorder.py dòng 122--138 --- decrypt\_clip()}
\begin{lstlisting}[language=Python, firstnumber=122]
def decrypt_clip(data: bytes,
                 password: str) -> Tuple[List[np.ndarray], float]:
    if len(data) < 44:           # dong 127: 32+12=44 byte minimum
        raise ValueError("Clip too short / corrupted.")
    salt       = data[:32]       # dong 129: parse salt
    nonce      = data[32:44]     # dong 130: parse nonce
    ciphertext = data[44:]       # dong 131: phan con lai
    key = derive_clip_key(password, salt)  # dong 132: re-derive key
    try:
        plaintext = AESGCM(key).decrypt(   # dong 134: decrypt + verify tag
                        nonce, ciphertext, None)
    except Exception:
        raise ValueError(
            "Decryption failed -- wrong password or corrupted clip.")
    return _unpack_frames(plaintext)       # dong 138: frames, fps
\end{lstlisting}
\end{code}

\begin{say}
``\textbf{Dòng 127}: guard --- data tối thiểu 44 byte (32 salt + 12 nonce).
Clip quá ngắn là bị lỗi hoặc không phải file hợp lệ.

\textbf{Dòng 129--131}: unpack theo offset cứng.
Salt luôn ở byte 0--31, nonce ở 32--43, còn lại là ciphertext.
Không cần delimiter vì kích thước header cố định.

\textbf{Dòng 132}: re-derive key --- \textit{cùng password + cùng salt = cùng key}.
Đây là tính deterministic của PBKDF2.
Key không bao giờ được lưu trên đĩa --- tạo tạm thời, dùng xong xoá khỏi memory.

\textbf{Dòng 134}: \texttt{AESGCM(key).decrypt} --- đồng thời verify Authentication Tag.
Nếu sai password $\to$ key sai $\to$ tag không khớp $\to$ exception.
Nếu ai đó sửa file dù 1 bit $\to$ tag không khớp $\to$ exception.
Dòng 136--137: bắt exception, raise ValueError với message \textbf{không tiết lộ chi tiết}.
Không nói ``salt sai'' hay ``key sai'' --- chỉ nói ``failed'' chung.
Information disclosure prevention.''
\end{say}

% ============================================================
\section{10:30--12:30 --- ADMIN\_AUTH.PY --- PBKDF2 + TIMING ATTACK}
% ============================================================

\begin{timing}{10:30 -- 12:30}
Chuyển tab sang \fpath{src/admin\_auth.py}.
\end{timing}

\begin{doact}
Click tab \fpath{admin\_auth.py}.
Nhấn \texttt{Ctrl+G}, gõ \textbf{31}, Enter.
\end{doact}

\begin{say}
``File cuối em muốn trình bày: \fpath{admin\_auth.py}.
Đây là xác thực admin --- bảo vệ quyền truy cập vào clip mã hoá và pipeline.''
\end{say}

\subsection{10:30--11:30 --- \_ITERATIONS và \_pbkdf2 (dòng 31--45)}

\begin{code}{admin\_auth.py dòng 31--45 --- \_ITERATIONS \& \_pbkdf2()}
\begin{lstlisting}[language=Python, firstnumber=31]
_ITERATIONS = 390_000   # OWASP 2023 minimum for PBKDF2-SHA256
_KEY_LEN     = 32       # 256-bit output key

def _pbkdf2(password: str, salt: bytes,
            iterations: int = _ITERATIONS) -> bytes:
    kdf = PBKDF2HMAC(
        algorithm=hashes.SHA256(),
        length=_KEY_LEN,        # 32 bytes = 256 bit
        salt=salt,
        iterations=iterations,  # dong 41: 390.000 lan SHA-256
    )
    return kdf.derive(password.encode("utf-8"))  # dong 43
\end{lstlisting}
\end{code}

\begin{say}
``\textbf{Dòng 31}: \texttt{\_ITERATIONS = 390\_000}.
Con số này \textbf{không phải em tự nghĩ ra}.
Đây là chuẩn tối thiểu của OWASP (Open Web Application Security Project) năm 2023
cho PBKDF2-SHA256.

Tại sao cần nhiều iterations đến vậy?
Hãy tưởng tượng attacker brute-force password 8 ký tự mixed-case + số.
Số khả năng: $96^8 \approx 7 \times 10^{15}$.
Nếu SHA-256 trực tiếp: máy hiện đại tính $10^9$ hash/giây $\to$ crack trong vài giờ.
Với PBKDF2 390.000 iterations: mỗi verify mất $390.000 / 10^9 = 0.39$ giây.
Brute-force cùng $7 \times 10^{15}$ khả năng: $7 \times 10^{15} \times 0.39s \approx$ hàng triệu năm.

\textbf{Dòng 41}: số iterations truyền vào PBKDF2HMAC.
\textbf{Dòng 43}: \texttt{kdf.derive} chạy SHA-256 đúng 390.000 lần và trả về 32 byte key.''
\end{say}

\subsection{11:30--12:00 --- \_ct\_equal và timing attack (dòng 48--50)}

\begin{doact}
Nhấn \texttt{Ctrl+G}, gõ \textbf{48}, Enter.
\end{doact}

\begin{code}{admin\_auth.py dòng 48--50 --- \_ct\_equal(): constant-time comparison}
\begin{lstlisting}[language=Python, firstnumber=48]
def _ct_equal(a: bytes, b: bytes) -> bool:
    """Constant-time comparison to prevent timing attacks."""
    return hmac.compare_digest(a, b)  # dong 50
\end{lstlisting}
\end{code}

\begin{say}
``Hàm nhỏ nhưng \textbf{rất quan trọng về bảo mật}.

Tại sao không dùng \texttt{a == b} thông thường?

Python \texttt{==} so sánh byte từ đầu đến cuối và \textbf{dừng ngay khi gặp byte đầu tiên khác}.
Nghĩa là: so sánh thành công 5 byte đầu rồi fail $\to$ tốn nhiều thời gian hơn
so với fail ngay byte đầu tiên.

Hacker thực hiện \textbf{Timing Attack}:
gửi 1000 request với password sai khác nhau,
đo thời gian phản hồi của server với độ chính xác millisecond.
Response lâu hơn $\to$ có nhiều byte prefix đúng hơn.
Qua thống kê, từng byte password có thể bị suy ra từ timing.

\textbf{Dòng 50}: \texttt{hmac.compare\_digest} luôn chạy \textbf{hết toàn bộ} byte,
bất kể sai ở vị trí nào.
Thời gian phản hồi không phụ thuộc vào vị trí sai $\to$ không có thông tin rò rỉ.''
\end{say}

\subsection{12:00--12:30 --- verify\_admin (dòng 83--95)}

\begin{doact}
Nhấn \texttt{Ctrl+G}, gõ \textbf{83}, Enter.
\end{doact}

\begin{code}{admin\_auth.py dòng 83--95 --- verify\_admin(): full authentication flow}
\begin{lstlisting}[language=Python, firstnumber=83]
def verify_admin(password: str,
                 path: str = _DEFAULT_CRED_PATH) -> bool:
    p = Path(path)
    if not p.exists():
        return False
    try:
        data = json.loads(p.read_text(encoding="utf-8"))
        salt      = base64.b64decode(data["salt"])  # dong 91
        stored    = base64.b64decode(data["hash"])  # dong 92
        iters     = int(data.get("iterations",
                                 _ITERATIONS))      # dong 93
        candidate = _pbkdf2(password, salt, iters)  # dong 94
        return _ct_equal(candidate, stored)          # dong 95
    except Exception:
        return False   # dong 97: fail silently
\end{lstlisting}
\end{code}

\begin{say}
``\textbf{Dòng 91--92}: đọc salt và hash từ JSON.
Salt + hash đều lưu dạng base64, không phải plaintext.

\textbf{Dòng 94}: \texttt{\_pbkdf2(password, salt)} re-derive.
Cùng password + cùng salt $\to$ cùng kết quả.
Đây là tính deterministic của PBKDF2 --- không cần lưu key.

\textbf{Dòng 95}: so sánh constant-time qua hàm dòng 50 vừa xem.

\textbf{Dòng 97}: \texttt{except Exception: return False}.
Bất kỳ lỗi nào --- JSON sai, base64 sai, file không tồn tại --- đều trả False.
Không có exception detail lọt ra ngoài.
Nguyên tắc: \textit{fail silently, no information disclosure.}''
\end{say}

% ============================================================
\section{12:30--13:30 --- DEMO LIVE}
% ============================================================

\begin{timing}{12:30 -- 13:30}
Chuyển sang browser.
\end{timing}

\begin{doact}
Mở tab browser \texttt{http://localhost:8000/admin/dashboard.html}.
Đứng vào khung camera (hoặc chỉ mặt trong webcam preview).
\end{doact}

\begin{say}
``Em demo nhanh 3 phần code vừa trình bày.''
\end{say}

\begin{doact}
\textbf{Demo 1 --- 4 preset:}
\begin{enumerate}
  \item Start pipeline
  \item Chọn preset \textbf{fast} --- pixelate, chỉ vào kết quả, nói ``dòng 55''
  \item Chọn \textbf{balanced} --- headcloak
  \item Chọn \textbf{strong} --- obliterate (hàm dòng 172) --- chỉ rõ không nhận ra mặt
  \item Chọn \textbf{strict} --- headcloak cực mạnh
\end{enumerate}
\end{doact}

\begin{say}
\screen\ ``Đây là 4 preset từ PRESETS dict dòng 53--91.
Mode strong ở đây chính là hàm \texttt{\_obliterate\_features} dòng 172 em vừa giải thích.
Nhìn kết quả: không thể nhận ra được gì.''
\end{say}

\begin{doact}
\textbf{Demo 2 --- Decrypt clip:}
Mở tab \texttt{http://localhost:8000/admin/secure-archive.html}.
Click một clip \texttt{.enc}.
\end{doact}

\begin{say}
``Tab này là Secure Archive --- clip gốc đã mã hoá bằng hàm \texttt{encrypt\_clip} dòng 109.

Nhập đúng password $\to$ xem clip gốc rõ nét.'' \pause

\textit{[sau khi clip hiện ra]} ``Clip gốc hiện ra đây.
Dòng 134: \texttt{AESGCM(key).decrypt} vừa chạy, verify Authentication Tag thành công.''
\end{say}

\begin{doact}
Nhập sai password, click Decrypt lại.
\end{doact}

\begin{say}
``Sai password: error message từ dòng 136--137: \textit{Decryption failed}.
Message không nói sai ở đâu --- đây là information disclosure prevention.''
\end{say}

% ============================================================
\section{13:30--14:00 --- KẾT THÚC VIDEO}
% ============================================================

\begin{timing}{13:30 -- 14:00}
Nhìn camera, nói kết thúc từ tốn.
\end{timing}

\begin{cam}
Nhìn thẳng camera. Nói chậm, rõ.
\end{cam}

\begin{say}
``Tóm lại, phần em đóng góp trực tiếp:

\fpath{privacy.py} --- 7 mode blur, quan trọng nhất là \texttt{\_obliterate\_features} dòng 172 với 4 bước phá huỷ không thể đảo ngược.

\fpath{anonymizer\_backend.py} --- 4 preset từ dòng 53, cân bằng giữa tốc độ và bảo mật.

\fpath{clip\_recorder.py} --- AES-256-GCM dòng 109--119, per-clip random salt và nonce, integrity guarantee.

\fpath{admin\_auth.py} --- PBKDF2 390.000 iterations theo OWASP 2023, constant-time comparison chống timing attack.

Tất cả hướng về một mục tiêu: \textbf{Privacy by Design} --- bảo vệ người dùng theo tiêu chuẩn công nghiệp.

Cảm ơn thầy cô.''
\end{say}

% ============================================================
\section{PHỤ LỤC --- CÂU HỎI DỰ KIẾN}
% ============================================================

\begin{tcolorbox}[colback=gray!5, colframe=BrickRed, breakable,
  title={\textbf{Q: Tại sao AES-GCM mà không phải AES-CBC?}}]
``GCM = Galois/Counter Mode --- authenticated encryption.
Ngoài mã hoá, GCM tạo Authentication Tag 128-bit.
Khi decrypt dòng 134, tag được verify đồng thời.
File bị sửa 1 bit $\to$ tag sai $\to$ fail ngay.
CBC chỉ mã hoá, không verify integrity.
Với bằng chứng pháp lý: integrity là bắt buộc, không thể dùng CBC.''
\end{tcolorbox}

\begin{tcolorbox}[colback=gray!5, colframe=BrickRed, breakable,
  title={\textbf{Q: 390.000 có đủ? Sao không dùng bcrypt hay Argon2?}}]
``390.000 là ngưỡng tối thiểu OWASP 2023 cho PBKDF2-SHA256 --- đủ cho use case này.
bcrypt và Argon2 tốt hơn về memory-hardness (chống GPU crack),
nhưng Python \texttt{cryptography} library đã có PBKDF2HMAC sẵn --- không cần thêm dependency.
Với đề án này, PBKDF2 390.000 đã đủ mạnh.
Production thực tế em sẽ cân nhắc Argon2id.''
\end{tcolorbox}

\begin{tcolorbox}[colback=gray!5, colframe=BrickRed, breakable,
  title={\textbf{Q: \_obliterate có thực sự không thể phục hồi không?}}]
``4 bước kết hợp: Nén 6\% + Scramble block + Blur 35px + Quantize 8 màu.
Thử nghiệm: chạy qua InsightFace buffalo\_l (model nhận diện tốt nhất em có access),
trả về similarity score $< 0.1$ với embedding gốc, ngưỡng nhận diện là 0.35.
Không nhận ra.
Về lý thuyết: bước nén 6\% đã irreversible --- thông tin bị mất permanently, không có super-resolution nào phục hồi được vì dữ liệu không còn tồn tại.''
\end{tcolorbox}

\end{document}
